This section contains an overview of already existing tools. This will help us
to concentrate on essential points and may serve as good consideration points
when we will be working on the design of our solution. 
%TODO:
There will also be provided a comparison table at the end of this section, which
will summarize the main properties of each solution.
% TODO:
% Properties to compare:
% - Open Source / Closed closed-source
% - Local \gls{llm} support / Remote LLM only
% - Context Awareness / No Context Awareness
% - Agentic Architecture / No Agentic Architecture
% - Shell Replacement / CLI Utility / Terminal Emulator

\subsection{Warp Terminal}

This is a developement-oriented and general-purpose \gls{ai}-enabled terminal
emulator, which uses remote Agent \gls{llm} System to enchance the experience of
users. The most of its code remains closed-source. It is not adapted for
executing on small \glspl{llm} and it provides no support for local \glspl{llm}.
Additionally, it is not a \gls{cli}/\gls{tui} application, but a full graphical terminal
emulator.
\\
This solution is popular and comercially successful, but it does not
completely fit into the scope of our work, as it is not designed to work with
local \glspl{llm}. However, it is worth mentioning, as it provides good AI-based
services for terminal users\cite{warp-web,warp-docs}.\\

\subsection{gptme}

This is general-purpose \gls{cli} utility with breadth \gls{ai} functionality. It is
compatible with OpenAI \gls{api}, which allows it to use both remote and local \glspl{llm}.
It uses Agentic architecture and implements \gls{mcp} protocol. This enables it to
use various tools (terminal, browser, code execution, etc.). It is also
context-aware, which means it can read and understand the files in the current
working directory and acknowledge host \gls{os}\cite{gptme-web,gptme-docs}.\\

Below the short interaction example with gptme is shown: % TODO inline code
\begin{lstlisting}[,
  caption={Example of gptme interaction with browser tool},
  label={lst:gptme},
  language=bash,
  style=code-listing,
  float
]
\$ gptme --tools=browser "How many calories are there in one spoon of sugar?"
[11:07:21] WARNING Skipping all confirmation prompts. cli.py:186
[11:07:21] Using browser tool with lynx
Found openai API key, using openai provider
Using model: openai/gpt-5
[11:07:22] Using logdir: ~/.local/share/gptme/logs/2025-12-10-running-sleepy-robot
Using workspace: ~/Workspace/bachelor/project/document
Skipped 2 hidden system messages, show with --show-hidden
--- ^^^ past messages ^^^ ---
User: How many calories are there in one spoon of sugar?
Assistant: Short answer: about 16 calories per level teaspoon.

Details:
- Level teaspoon (≈4 g sugar): ~15–16 kcal
- Heaped/rounded teaspoon (≈5–7 g): ~20–28 kcal
- Tablespoon (≈12–15 g): ~48–60 kcal

Sugar has ~4 kcal per gram (sucrose ~3.87 kcal/g), so you can estimate by weight if needed.
\end{lstlisting}

\begin{itemize}
  \item Uses OpenAI API for accessing remote/local \glspl{llm}
  \item Uses Agent architecture and MCP tools
  \item Is context-aware and is able to use tools (terminal, browser,
    code execution, etc.)
  \item Implements MCP protocol
\end{itemize}
\\


% TODO: citations for models

When tested manually with local \glspl{llm} (Gemma3 4b, Qwen2.5 1.5b), it did not
perform well, as the models were not able to handle the complexity of the
tasks. It usually due to too long prompts which made \glspl{llm} halucinate on too
many details. This shows, that the solution is not fully adapted for local
\glspl{llm}, which is one of the main goals of our work. This should be considered
when designing our solution to use more fine-grained prompts for local \glspl{llm} and
optimize the architecture for local \glspl{llm}.\\

\subsection{gemini-cli}

% chat-like vs chat?
This a is general-purpose \gls{cli} chat-like application by Google. It uses Gemini
\glspl{llm} to provide \gls{ai} capabilities in terminal. It is useful for various tasks,
including code generation, data analysis, and general Q\&A. It also supports
\gls{mcp} protocol and custom tool integration. Hovewer, it is dependable on remote
\glspl{llm} and cannot be used with local \glspl{llm}\cite{gemini-web,gemini-docs}.\\

\subsection{Yai (Your \gls{ai} Assistent)}

This application is an interactive \gls{tui} chat application, which uses OpenAI to
provide user with \gls{ai} capabilities in terminal. It supports tools interaction
and basic agentic architecture. In current form, it is not able to use local
\glspl{llm}. Hovewer, it is possible to adapt it for local \gls{llm}s with some modifications
in its source code.\cite{yai-web,yai-docs}.\\

% TODO Worth testing later?

\subsection{Aider}

This is an \gls{llm}-powered code agent with focus on software development tasks and
terminal integration. It supports various tools and is context- and
project-aware. It is able to use both remote and local \glspl{llm}. It is also
open-source.\cite{aider-web,aider-docs}.\\

\subsection{termai}

This is a terminal-based \gls{llm} assistant similar to gptme, written in Rust. It
supports various \gls{llm} providers, but it is not able to use local \glspl{llm} in its
current form. It is open-source and has basic agentic architecture\cite{termai-web}. \\

\subsection{AiTerm}

The application that utilizes tmux terminal multiplexer to access the terminal
\gls{io}. It sends the terminal output as context to chosen \gls{llm} and gets the
response back. It has execution agent capabilities, which allows it to run
commands in terminal based on \gls{llm} response. It is open-source and is adapted to
use with OpenAI \gls{api}, thus and it can be used with local \glspl{llm} as well.\\

We can draw inspiration from its approach to Bash integration using tmux sessions to
provide seamless interaction between \gls{llm} and terminal. Hovewer, this approach
turned out to be somewhat clunky in practice, as it couldn't handle well the
~\gls{tui} applications, but only get bash's responses\cite{aiterm-web}. \\


\subsection{Other solutions}

There are also such solutions, like \textsl{ai-terminal-assistent (github)}{}
or \textsl{cmd-ai (github)}, which are just wrappers for Rest \gls{api} calls to
remote \glspl{llm}. They provide basic \gls{ai} capabilities in terminal, but
have no significant features.

% % comparison table
% % Split into two tables
%
% \begin{table}[ht]
%   \small
%   \begin{tabular}{|l|c|c|c|c|c|}
%     \toprule
%     \textbf{Property} & \textbf{Warp Terminal} & \textbf{gptme} &
%     \textbf{gemini-cli} & \textbf{AiTerm} & \textbf{Yai} \\
%     \midrule
%     \textbf{Open Source} & No & No & TBD & Yes & No \\
%     \textbf{Context Awareness} & Partial (Display Reading) & Yes &
%     TBD & Partial & No \\
%     \textbf{Agentic Architecture} & Yes & Yes & TBD & Yes & No \\
%     \midrule
%   \end{tabular}
%   \\
%   \small
%   \begin{tabular}{|l|c|c|c|c|}
%     \toprule
%     \textbf{Property} & \textbf{aider} &
%     \textbf{ai-terminal-assistant} & \textbf{cmd-ai} & \textbf{termai} \\
%     \midrule
%     \textbf{Open Source} & TBD & TBD & TBD & TBD \\
%     \textbf{Context Awareness} & TBD & TBD & TBD & TBD \\
%     \textbf{Agentic Architecture} & TBD & TBD & TBD & TBD \\
%     \midrule
%   \end{tabular}
%   \caption{Second half of comparison table for Local \gls{llm} Assistant Development}
% \end{table}
